\documentclass{article}

\usepackage[natbib=true, style=apa, alldates=year]{biblatex}
\usepackage[pdfauthor={Julian Barg},pdftitle={Research Summary}]{hyperref}
\usepackage{lipsum}
\usepackage{rotating}
\usepackage[notablist,tabhead,nofiglist,fighead]{endfloat}
\usepackage{nameref}
\usepackage{tabularx}
\usepackage{graphicx}
\usepackage{tikz}
\usepackage{booktabs}

\usetikzlibrary{decorations.pathmorphing}

\addbibresource{bibliography.bib}

\hyphenation{Trans-Canada}

\title{Research Summary}

\author{
	Barg, Julian\\
	\texttt{jbarg.phd@ivey.ca}\\
	Ivey Business School
}

%!TEX root = ../main.tex

\newcommand{\timeline}{\hspace{-2.3pt}$\bullet$ \hspace{5pt}}

\newcommand{\rota}[1]{\rotatebox{65}{#1}}

\newcommand{\ye}{
	\begin{tikzpicture}
		\draw[thick] (0,0) circle (1.0mm);
		\fill (0,0) circle (0.4mm);
	\end{tikzpicture}
}

\newcommand{\no}{\tikz \draw[thick,color=black!50] (0,0) circle (1.0mm);}

\newcommand{\fu}{
	\begin{tikzpicture}
		\draw[->,thick,decorate,decoration={snake,amplitude=0.4mm,segment length=1mm,pre length=0.1mm,post length=0.5mm}] (0,0) -- (0,-0.26);
		\draw[rounded corners=1pt] (0,0.08) -- (-0.23,-0.28) -- (0.23,-0.28) -- cycle; 
	\end{tikzpicture}
}

\newcommand{\hm}{
	\begin{tikzpicture}
		\draw[thick] (0,.6mm) arc(90:270:1.0mm);
		\fill (0,0) arc (90:270:0.4mm);
	\end{tikzpicture}
}

\newcolumntype{s}{p{0.3cm}}

\newcolumntype{P}[1]{>{\raggedright\arraybackslash}p{#1}}

\DeclareDelayedFloatFlavor{sidewaystable}{table}

\sloppy

\begin{document}

	\maketitle
	\clearpage

	%!TEX root = ../main.tex

What is sustainable? What is sustainability? There are some formulaic definitions that are widely accepted, including the dogmatic definition that "[s]ustainable development is development that meets the needs of the present without compromising the ability of future generations to meet their own needs" \citep[IV]{WCED1987}. But if you ask two sustainability scholars to define sustainability in their own terms, chances are you are still going to get very different answers. Potentially including the ubiquitous phrase "I know it when I see it" \citep{White2013}. All that is not to criticize sustainability scholarship. It is the nature of the beast of dealing with complex systems--such as our planet--that things are difficult to nail down.

Definitions attain importance in research because they shape the research that takes place after the introduction section. Things may be assumed away, and a definition can determine the direction of a research project. The definition of sustainable development by the World Commission on Environment and Development (WCED)--despite being the most widely used definition of sustainability by far--is still contested. On the one hand, the WCED definition is widely used as a working definition. On the other hand, the WCED definition is criticized for putting too much focus on economic \textit{development} for resolving current and future distributive and environmental problems \citep[e.g.,][6f.]{Constanza2014a}.

This blurb introduces the discourse on the WCED definition and what assumptions it introduces to research. It covers on the one hand research in business journals that apply the definition. These articles have in common that they turn to businesses and their policy decisions for potential solutions to common environmental or social problems. On the other hand, there is a critical 
% \todo{Avoid label "critical" for its connotations?} 
discourse on sustainable development that explicitly discusses the inherent assumptions and typically rejects the notions that companies are suitable agents that can be relied on for change \citep[e.g.,][]{Banerjee2003}.
% This critical  discourse maintains that action by businesses is and will remain insufficient, and our current inaction outside the business realm will be costly in environmental and economic terms.

The discourse is driven by competing assumptions on the ability for corporations to act as agentic actors for change. 
\begin{itemize}
	\item \hl{Add "third way" implications similar to \citet{Giddens1979}}.	
	\item Actually in some sense, this directly parallels Gidden's comments on structure vs. agency--e.g., \citet{Heikkurinen2019}.
	\item the "third way" is actually the extant literature I have extracted before--\url{http://wiki.jbarg.net/Constructivism%20%26%20Ecology.md#environmental-science-stream}
\end{itemize}


	%!TEX root = ../main.tex

\subsection*{Studies}

This section introduces three studies that are designed to generate research insights from the frail progress on the deinstitutionalization of oil and gas that is exemplified by the revocation of the permit for Keystone XL. The remainder of this document, after a brief summary of the \textit{quantitative} data I have collected on the industry, focuses on the \textit{qualitative} methods and data for the second study. The studies draw on institutional theory and organizational learning. Institutional theory is useful for explaining the stability of the current structure, and the challenges around substantial (as opposed to symbolic) change \citep{Meyer1977,Hoffman2015}. Organizational learning on the other hand emphasizes the potential for adaption, knowledge sharing, and vicarious observation, while keeping one eye on obstacles \citep{Madsen2010,Madsen2018,Rerup2021}.

Underlying all three studies is the understanding I have briefly developed above: that the polarization of the discourse on oil and gas has not been overcome by the defeat of Keystone XL. That the recent institutional developments in this area are characterized by twists and turns. And most importantly, that we should pay more attention to pipeline projects such as Keystone XL as potential precursor of the deinstitutionalization--or resurrection!-- of oil and gas, and "expect the unexpected".

The \textit{first study} focuses on the industry side of the discourse. The operation and construction of pipelines requires a social license. The US experiences over 100 significant
\footnote{The Pipeline and Hazardous Materials Safety Administration describes a pipeline spill as significant if the spill volume exceeds 50 barrels, or if one of three other criteria is met--see \url{https://www.phmsa.dot.gov/data-and-statistics/pipeline/pipeline-incident-flagged-files}, accessed 2021-05-13.} 
pipeline spills every year.
\footnote{See \url{https://www.jbarg.net/incident_dashboard.html}, accessed 2021-05-13.} Pipeline operators routinely make references to technology--both clean-up technology, and pipeline safety technology--in their statements on recent pipeline spills. Since both the quality of implementation and effectiveness of these new safety technologies is highly opaque for external audiences, the risk of decoupling and symbolic adaption is high \citep{Wijen2014a}. In other words, pipeline operators can use symbolic references to technology in conjunction with selective disclosure of information on spill causes as a means of greenwashing \citep{Lyon2015}. Study one explores the potential of technology as a device for greenwashing, using a mix of quantitative data on pipeline networks and spills, and qualitative descriptions of spill causes and responses.

The first study's theoretical focus on greenwashing limits the analysis to only one side of a richer discourse. Five communities have participated in the discourse on Keystone XL: pipeline operators, environmental activists, political decision makers, local residents, and the scientific community. The \textit{second study} opens up the blinds to provide the complete picture of this multivocal environment. Participants meet in public forums across multiple years \citep{Latour2005} and produce their own texts before and after, allowing us to repeatedly observe their beliefs, information exchange, and strategizing over time. \citet{Maguire2009} describe deinstitutionalization as a process of discourse and collective alteration of institutions. In the case of Keystone XL, we observe only a partial reconfiguration of knowledge networks, which leads to a frail deinstitutionalization of large oil and gas infrastructure projects \citep[cf.][]{Latour1984}. Study two studies focuses on the different communities that participate in the discourse on Keystone XL \citep{Aronczyk2019,KnorrCetina1999}, and how their participation on the Keystone XL discourse affects their understanding of the environmental risks and impacts related to pipelines and fossil fuels. Finally, I develop a model of the social process of deinstitutionalization in a multivocal environment.
% The ultimate rejection of the Keystone XL application neither resolves the conflicts nor does it lead to a stable new status quo. The conflict has given rise to new networks of pipeline opponents. But the sparring communities and their fundamental disagreements persists after the conflict. The temporary new status quo can be attributed to the new constellation that has emerged between the belligerents \citep{Latour1984}.

The \textit{third study} is a theory paper that applies the insights from studies one and two to the theory on organizational and collective learning. Whether a pipeline operators can maintain their public license does not only depend on the accuracy of the representations of their assets that they share with their audience. Rather, their models of the world need to become shared across a sufficient share of audience members. The learning literature similarly emphasizes the social dimension of learning--think vicarious observation \citep{Madsen2018}. The literature distinguishes between two dimensions--reliablity and validity \citep{March1991a}. Reliability describes how widely shared a model of the world has become. Validity describes how useful the model is for prediction and control \citep{Rerup2021}. Contrary to our intuition, models that are suboptimal in terms of validity can become widely shared in the world \citep{Levitt1988}. The third chapter is dedicated to the theory of reliability and validity.

	%!TEX root = ../main.tex

\subsection*{Data on the pipeline industry}

The studies are informed by three datasets on the pipeline industry that I have collected in preparation of the dissertation research. The Pipeline and Hazardous Materials Safety Administration (PHMSA) is the preeminent source of quantitative data on the pipeline industry in the US. PHMSA provides data on pipeline miles and the volume of oil transported (in gallon-miles) by operator every year. PHMSA also maintains a repository on every pipeline spill in the US. This repository includes the data we need to carry out study one--the location, volume, cause, and a narrative of every of the over 100 significant spills (>50 gallon, or meeting one of three other conditions) that occur in the US every year.

In addition, the National Transportation Safety Board (NTSB) provides in-depth narratives of a small number of noteworthy pipeline incidents over an extended period of time. The NTSB uses its own resources to investigates these incidents and provides in-depth reports of usually 10-70 pages as well as shorter briefs. Since 1996, the NTSB has completed 53 investigations on pipeline incidents. These incidents provide an understanding of both the organizational and technological causes of pipeline spills \citep[cf.][]{Perrow1984}.

	%!TEX root = ../main.tex

\section*{Data Collection}\label{sec:data}

% Data collection for study two involves five steps. First, a timeline of events is created \citep{Maguire2009}. From that timeline, we select forums that facilitate a discourse \citep{Latour2005}. Then, texts from the participating communities on the discourse are gathered. The data analysis is carried out on the basis of these texts. Finally, I interview individual members of the communities to test the robustness of my fiends and ensure convergence.

This section introduces the data collection for study two. To understand the discourse on Keystone XL, and how it unfolded over time, we require a sample of texts that encompasses multiple relevant actors over an extended period of time. The first step toward building this corpus is to develop a timeline of events \citep{Maguire2009}. The basis for that timeline are newspaper articles from the sampling period of 2008-2021. These articles were obtained from five major American national news outlets
% The Wall Street Journal, The New York Times, Los Angeles Times, and The Washington Post 
through Factiva by searching for the Keywords "Keystone XL", and "KXL". Further, articles on Keystone XL over the sampling period were collected from the Oil \& Gas Journal. 
% An excerpt of the timeline can be seen in \autoref{tab:timeline}.

% %!TEX root = ../main.tex

\begin{table}
\caption{Keystone XL--Timeline of selected events}
\label{tab:timeline}
\begin{tabularx}{\textwidth}{r @{\hspace{0.5\tabcolsep}} l |@{\timeline} X}
% \begin{tabular}{r @{\hspace{0.5\tabcolsep}} l |@{\timeline} >{\raggedright\arraybackslash}p{9cm}}
\toprule

% 2005 & Feb 	 & TransCanada unveils plans to build the Keystone pipeline to transport diluted bitumen from Alberta, Canada to Texas.\\
2008 & Mar   & The Bush administration issues the Presidential Permit for the \textit{Keystone Pipeline}, predecessor of the \textit{Keystone XL} pipeline.\\
	 & June  & TransCanada proposes the Keystone XL pipeline.\\
2010 & April & The Deepwater Horizon oil platform explodes before spilling 4.9 million barrels of oil in the the Gulf of Mexico.\\
	 & June  & Phase 1 of Keystone begins delivering oil from Alberta, Canada to Steele City, Nebraska in the US.\\
	 & July  & An Enbridge pipeline spills over 20,000 barrels of diluted bitumen into the Kalamazoo River in Michigan.\\
2012 & Jan   & The Department of State rejects the Keystone XL application after congress passes legislation to require the completion of the National Interest Determination within 60 days.\\
2014 & Jan	 & The Keystone pipeline is completed and begins delivering oil from Canada to Texas.\\
2015 & Nov	 & The Department of State rejects the second Keystone XL application.\\
2016 & Jan 	 & In his first week in office, Trump signs a presidential memorandum that invites TransCanada to resubmit its application.\\
2017 & Mar   & Trump signs the Presidential Permit, based on a slightly revised National Interest Determination that now rules in favor of Keystone XL. The permit later faces challenges in court on procedural grounds.\\
2019 & Mar 	 & Trump revokes the permit for Keystone XL and issues a new permit. Legal challenges continue.\\
2020 & April & Construction starts on Keystone XL.\\
2021 & Jan 	 & Biden revokes the permit for Keystone XL.\\
\bottomrule

\end{tabularx}
% \end{tabular}
\end{table}

The timeline on Keystone XL includes multiple public forums \citep{Latour2005}. These public forums serve two functions in this study. (1) They allow for the identification of relevant participants in the discourse. 12 types of participants participate in the forums, 6 of which are specific organizations that participate repeatedly. 
% (see \autoref{tab:forums}). 
(2) The public forums are a compressed version of the discourse itself, where actors lay out their understanding of the permit application process, of Keystone XL, and of its environmental impact. The actors also respond to each others claims in this forum, laying bare the extent to which they have engaged with each others' arguments. I.e., a rebuttal testimony begins with a summary of the other side's position, and make a useful gauge of engagement and comprehension.

% %!TEX root = ../main.tex

\begin{sidewaystable}

\tiny
\centering\caption{Forums on Keystone XL}
\label{tab:forums}
\begin{tabular}{p{3.5cm}@{\hskip 5mm} r s s s s s s s s s s s s s s s}
	\toprule

	Event & Time 
		& \rota{Keystone} 
		& \rota{State government} 
		& \rota{Federal government} 
		& \rota{Local indigenous organization(s)} 
		& \rota{Local residents}
		& \rota{Labor unions}
		& \rota{Dakota Rural Action} 
		& \rota{Bold Nebraska}
		& \rota{Sierra Club} 
		& \rota{350.org}
		& \rota{Indigenous Environmental Network}
		& \rota{TransCanada whistleblower}
		\\

	\midrule

	Canada National Energy Board--OH-1-2009 & Mostly 2009 & 
		\ye & \no & \ye & \ye & \hm & \ye & \no & \no & \ye & \no & \no & \no\\

	South Dakota Public Utilities Commission--HP09-001 & Mostly 2009 &
		\ye & \ye & \no & \no & \no & \no & \ye & \no & \no & \no & \no & \no\\

	Montana Department of Environmental Quality--Certificate of Compliance & March 2012 &
		\ye & \ye & \no & \no & \no & \no & \no & \no & \no & \no & \no & \no \\

	Senate of Canada--TransCanada whistleblower testimony & June 6, 2013 &
		\fu & \no & \ye & \no & \no & \no & \no & \no & \no & \no & \no & \ye\\

	South Dakota Public Utilities Commission--HP14-001 & Mostly 2014 & 
		\ye & \ye & \no & \ye & \ye & \no & \ye & \ye & \ye & \ye & \ye & \ye\\

	Nebraska Public Service Commission--OP-0003 & Mostly 2017 & 
		\ye & \ye & \no & \ye & \ye & \ye & \no & \ye & \ye & \no & \no & \no\\

	Montana Department of Environmental Quality--401 Water Quality Certification & Mostly 2020 &
		\ye & \ye & \no & \no & \ye & \no & \no & \ye & \no & \no & \no & \no\\

	\bottomrule
\end{tabular}

\end{sidewaystable}

By selecting on the participation in forums, we obtain a sample of actors that are both knowledgeable and engaged. To participate in the public forum, an organization needs to both familiarize itself with the formal process and strategize, as well as obtain and organize expert knowledge. 
% Some of the participants, such as the Sierra Club or Bold Nebraska are or become communications experts over the course of the discourse.

The final step toward creating a sample of texts is to gather the communication of the actors on Keystone XL over the time of the discourse. Most of the communities have formed and organize themselves through websites. Documents on TransCanada, Keystone XL, and the Sierra Club are also published online. The documents are downloaded by hand, or automatically through web scraping where appropriate. In some cases, the texts are automatically filtered to retain only those that discuss Keystone XL, using the keywords "Keystone XL" and "KXL". The filtered texts, and the texts from step two constitute the final sample that is imported into ATLAS.ti for the data analysis.

% Some of the forums were open to public participation. Comments by residents were also collected but were too numerous to manually sight. Topic modeling and sentiment analysis were used to identify substreams of that corpus. Samples from each substream then entered the analysis.

	%!TEX root = ../main.tex

\subsection*{Methods}

The purpose of this study is to extract the meaning of words and phrases from text \citep{Parker1992}. We identify the main themes of the discourse by reviewing and coding the public forums. To decode their meaning, we first reference the timeline of events that was compiled in the first step in the section \nameref{sec:data} \citep{Maguire2009}. We then turn to the texts that the community has produced ahead of the forum. Connecting themes from the forums with texts that were produced ahead of the forum, allows us to more accurately decode their meanings. Texts that are specific to the forum sometimes reveal a community's strategy for the forum.

We have now obtained coded meanings and strategies on the forum. The next step is to compare coded themes before, during, and after. To obtain a process model of the discourse, we code the texts produced by all communities ahead and after each forum. Discrepancies between the before and during provide more information on the communities' strategies. Discrepancies between the during and after suggest an evolving understanding through exchange at the forum. Finally, we track themes across communities and across time throughout multiple forums. This step allows for the identification of more connections and overarching trends.

In addition to the analysis of texts, this analysis uses interview data. The primary function of the interviews is to test the robustness of our findings. For that purpose, I compile timelines on the discourse for each community. These timelines are shared with interviewees from the communities to triangulate events and ensure completeness of the data. Where data is incomplete, we return to the previous step and review the findings, before carrying out more interviews.


	% The first study analyses the rhetorical strategies of the pipeline industry from a \textit{greenwashing} perspective \citep{Lyon2015}. The pipeline industry has a generous social license that--hundreds of individual companies are allowed operate facilities that carry enormous risks for the local environment. The US experiences over 100 significant\footnote{The Pipeline and Hazardous Materials Safety Administration describes a pipeline spill as significant if the spill volume exceeds 50 barrels, or one of three other criteria is met--see \url{https://www.phmsa.dot.gov/data-and-statistics/pipeline/pipeline-incident-flagged-files}, accessed 2021-05-12.} pipeline spills every year.\footnote{See \url{https://www.jbarg.net/incident_dashboard.html}, accessed 2021-05-12.} To maintain its social license, and expand its footprint, the industry has to navigate a contradiction. The industry has to present a technology with obvious risks as perfectly controllable and safe. How does the industry, and individual operators navigate this rhetorical challenge?

	\clearpage
	\printbibliography

\end{document}